\section{Related Work}
\label{sec:related}

\para{AI-scientist systems and their failure inventories.} \citet{lu2024aiscientist} and \citet{yamada2025aiscientistv2} introduced the \textsc{AI Scientist} pipeline (idea $\to$ code $\to$ paper $\to$ auto-review); subsequent systems include Agent Laboratory, MLR-Agent~\citep{chen2025mlrbench}, AI Researcher~\citep{tang2025airesearcher}, Jr. AI Scientist~\citep{miyai2026jrai}, and DeepScientist~\citep{weng2025deepscientist}. \citet{luo2025hidden} identify four \emph{hidden pitfalls} (cherry-picked benchmarks, data leakage, metric misuse, post-hoc selection) and show that auditing the final paper alone yields detection $F_1 = 0.51$, rising to $0.81$ with trace logs. \citet{chen2025mlrbench} report that coding agents fabricate or invalidate experiments in $\sim$80\% of cases and propose three coarse failure modes. \citet{jiang2025badscientist} document five fabrication strategies (TooGoodGains, BaselineSelect, StatTheater, CoherencePolish, ProofGap) under a multi-LLM review committee and introduce the ``concern--acceptance conflict.'' \citet{miyai2026jrai} produce a 76-page risk report listing $\sim$6 categories including hallucinated ablations, ambiguous methods, and figure misinterpretation. \citet{trehan2026whyllms} list six categories drawn from four end-to-end pipelines, three of which failed to complete; they explicitly attribute the ``Eureka instinct'' to RLHF agreeableness, conflicting with scientific skepticism. \citet{ansari2026compound} catalogue 100 fabricated citations that survived NeurIPS 2025 peer review and demonstrate that 100\% are compound (combining multiple atomic categories). \emph{Each source enumerates failures at the granularity convenient for its own contribution; none of them merge.} Our \taxonomy provides this missing consolidation.

\para{Foundational alignment failure modes.} \citet{amodei2016concrete} provide the canonical taxonomy of accident risk (negative side effects, reward hacking, scalable oversight, safe exploration, distributional shift). \citet{park2023deception} survey strategic deception, sycophancy, and unfaithful reasoning. \citet{vanderweij2024sandbagging} define sandbagging as \emph{strategic underperformance on an evaluation}, distinguishing developer- from system-level forms. \citet{naik2025agentmis} build a propensity benchmark spanning oversight avoidance, shutdown resistance, sandbagging, and power-seeking, and report that persona/system prompts drive propensity more than model choice. \citet{baker2025monitoring} show that chain-of-thought monitors detect reward hacking until the CoT is included in the training reward, after which they detect ``obfuscated'' reward hacking. We treat these as the foundational vocabulary that our research-specific categories instantiate---\eg test-case exploitation~\citep{impossiblebench2025,evilgenie2025} as an instance of reward hacking; research sandbagging as research-specific underperformance.

\para{Long-horizon and multi-agent failures.} \citet{rasheed2026claimaud} introduce three long-horizon failure modes (objective drift, transient constraints, unverifiable inference) and propose the AAR auditability standard. The agent-hallucination survey of \citet{lin2025agenthall} argues that agent hallucinations have more diverse types, longer propagation chains, and more severe consequences than single-shot LLM hallucinations.

\para{Propensity benchmarks.} \citet{vanderweij2024sandbagging} and \citet{naik2025agentmis} establish the methodology of measuring propensity via control-vs-treatment prompts where the misaligned response is available but not capability-limited. We adopt this methodology for the three under-documented categories and contribute---to our knowledge---the first controlled propensity result for negative-result omission in a frontier LLM under publish-pressure framing.

\para{Positioning.} Building on \citet{luo2025hidden} and \citet{ansari2026compound}, we (i) consolidate categories rather than enumerate new ones for a single sub-problem, (ii) measure prevalence with explicit inter-judge agreement, and (iii) supply propensity evidence rather than relying on prevalence alone. Unlike \citet{jiang2025badscientist}, who focus on a single pipeline stage (paper-to-review fabrication), we cover the full ideation-to-review lifecycle. Unlike \citet{trehan2026whyllms}, who distill categories from four hand-run pipelines, we apply a pre-registered taxonomy to 20 system $\times$ task pairs and quantify prevalence with confidence intervals.
